\chapter{Discrete Dipole Approximation}

In this appendix we present a detailed description of the Discrete Dipole Approximation implemented during the thesis.
This method, initially introduced by DeVoe \cite{devoe1964optical} to study the optical properties of molecular aggregates, has been improved over the years. A significant contribution was made by Purcell and Pennypacker \cite{Purcell_1973}, by including retardation effects.\\
In this thesis, the studied disordered systems where formed by particles positioned at some point in space. Each particle, described by its polarizability $\boldsymbol{\alpha}$, behaves like a dipole when under the influence of an electromagnetic field. This induce dipole is described by:

\begin{equation}
mathbf{p}_i = \epsilon_{0}\alpha\left(\omega\right)\mathbf{E}_i
\end{equation}

where $\mathbf{E}_i$ is the polarizability of the particle.


The consequence of the particle's excitation is the 
